\chapter{Measurement Data and Reality}

\section{Motivation}

\begin{remark}
In practical applications, alignment geometry is rarely available in a
clean, analytic form. Instead, it is derived from heterogeneous data
sources with varying levels of precision, structure, and semantic
completeness.
\end{remark}

\begin{remark}
Therefore, the modelling process must explicitly distinguish between
\begin{itemize}
\item observed data,
\item interpreted structure,
\item and reconstructed geometry.
\end{itemize}
\end{remark}

---

\section{Sources of Data}

\begin{definition}[Data source]
A data source is any origin of geometric or alignment-related
information.
\end{definition}

\begin{remark}
Typical sources include:
\begin{itemize}
\item survey measurements,
\item legacy engineering data,
\item container formats such as \LandXML,
\item railway-specific formats,
\item point clouds and derived datasets.
\end{itemize}
\end{remark}

---

\section{Nature of Real-World Data}

\begin{remark}
Real-world data exhibits several characteristic properties:
\end{remark}

\begin{itemize}
\item \textbf{Incompleteness:} Not all required information is present.
\item \textbf{Noise:} Measurements are subject to error.
\item \textbf{Ambiguity:} Multiple interpretations may be possible.
\item \textbf{Inconsistency:} Different parts of the data may contradict each other.
\end{itemize}

\begin{remark}
As a consequence, data cannot be directly equated with geometry.
\end{remark}

---

\section{Coordinate Systems}

\begin{remark}
Geometric data is always embedded in a coordinate reference system (CRS).
\end{remark}

\begin{definition}[Coordinate reference system]
A coordinate reference system defines how numerical coordinates relate
to positions in physical space.
\end{definition}

\begin{remark}
In infrastructure applications, multiple systems may coexist, including:
\begin{itemize}
\item global systems (e.g. WGS84),
\item projected systems (e.g. UTM, Gauss--Krüger),
\item domain-specific systems (e.g. railway reference systems).
\end{itemize}
\end{remark}

\OPEN{Formal integration of CRS handling into the alignment model.}

\RESEARCH{Transformation consistency across heterogeneous data sources.}

---

\section{Interpretation Layer}

\begin{definition}[Interpretation]
Interpretation is the process of assigning structure and meaning to raw data.
\end{definition}

\begin{remark}
Typical interpretation tasks include:
\begin{itemize}
\item identifying alignment segments,
\item reconstructing curvature behaviour,
\item associating auxiliary data such as profile or cant,
\item resolving ambiguities in topology.
\end{itemize}
\end{remark}

\begin{remark}
Interpretation is inherently non-unique and often requires domain
knowledge.
\end{remark}

\OPEN{Formal separation between data and interpretation.}

---

\section{From Data to Sparse Model}

\begin{remark}
The sparse alignment model provides a target structure for data
interpretation.
\end{remark}

\begin{definition}[Reconstruction from data]
Reconstruction is the process of deriving a sparse alignment model from
a given data source.
\end{definition}

\begin{remark}
This process typically involves:
\begin{itemize}
\item segmentation of data,
\item classification of element types,
\item estimation of curvature laws,
\item assignment of initial conditions.
\end{itemize}
\end{remark}

\TODO{Formal pipeline from raw data to sparse alignment.}

\OPEN{Handling of incomplete or conflicting input data.}

---

\section{Container Formats}

\begin{remark}
Many real-world datasets are provided in structured container formats.
\end{remark}

\begin{remark}
Examples include:
\begin{itemize}
\item \LandXML,
\item railway-specific XML formats,
\item spreadsheet-based engineering datasets.
\end{itemize}
\end{remark}

\begin{remark}
Such formats often mix geometric, semantic, and metadata in ways that
are not directly aligned with the sparse model.
\end{remark}

\RESEARCH{Mapping of container formats to sparse representation.}

\TODO{Definition of intermediate formats (e.g. landFAT).}

---

\section{Uncertainty and Error}

\begin{remark}
All real-world data is subject to uncertainty.
\end{remark}

\begin{definition}[Measurement uncertainty]
Measurement uncertainty describes the deviation between observed data
and the true physical quantity.
\end{definition}

\begin{remark}
Uncertainty propagates through:
\begin{itemize}
\item interpretation,
\item reconstruction,
\item numerical integration.
\end{itemize}
\end{remark}

\RESEARCH{Propagation of uncertainty through the reconstruction pipeline.}

\OPEN{Integration of uncertainty into the sparse model.}

---

\section{Topological vs Geometric Information}

\begin{remark}
Data may contain both geometric and topological information.
\end{remark}

\begin{definition}[Geometric information]
Geometric information describes positions, directions, and curvature.
\end{definition}

\begin{definition}[Topological information]
Topological information describes connectivity and relationships between elements.
\end{definition}

\begin{remark}
In many datasets, these two aspects are interwoven and not cleanly separated.
\end{remark}

\OPEN{Separation and integration of topology and geometry in the model.}

---

\section{Engineering Reality}

\begin{remark}
Engineering data is often shaped by practical constraints rather than
mathematical clarity.
\end{remark}

\begin{remark}
Typical phenomena include:
\begin{itemize}
\item piecewise definitions without explicit continuity guarantees,
\item implicit assumptions,
\item domain-specific conventions,
\item historical artefacts in legacy data.
\end{itemize}
\end{remark}

\RESEARCH{Systematic classification of real-world alignment data quality.}

---

\section{Role within the Overall Framework}

\begin{remark}
The data layer connects the abstract modelling framework to real-world
applications.
\end{remark}

\begin{remark}
It forms the entry point for:
\begin{itemize}
\item import pipelines,
\item validation procedures,
\item reconstruction workflows.
\end{itemize}
\end{remark}

---

\section{Summary}

\begin{summary}
Real-world alignment data is heterogeneous, noisy, and often ambiguous.

A clear separation between raw data, interpretation, and reconstructed
geometry is essential.

The sparse alignment model provides a structured target for
reconstruction, while numerical methods and domain knowledge are
required to bridge the gap between data and geometry.
\end{summary}
